% Figure 2 — Pearl's Causal Ladder for FactoryBench
% Layout: vertical ladder, L1 (bottom) to L4 (top), each row is a horizontal bar
% with: icon | level label + formula | example question
% Color intensity increases from L1tint (lightest) to L4tint (darkest)

\begin{figure}[t]
    \centering
    \begin{tikzpicture}[
        every node/.style={font=\sffamily},
        row/.style={
            minimum width=\linewidth-4pt,
            minimum height=1.42cm,
            inner sep=0pt,
            outer sep=0pt,
        },
    ]

    % Vertical spacing between rows
    \def\rowsep{1.52}

    % --- Row geometry ---
    % Each row: tinted left accent bar | icon | title + formula | example
    % We draw bottom-up so the ladder ascends visually

    \foreach \i/\tint/\icon/\lbl/\subtitle/\formula/\example in {%
        0/L1tint/{\faEye}/L1/Association/{$P(E \mid S,\, C)$}/{``When the payload is heavier, the motor current is higher.''},%
        1/L2tint/{\faCog}/L2/Intervention/{$P(E \mid \mathrm{do}(S))$}/{``If I command $+$20\% torque now, the current rises to ${\sim}X$\,A.''},%
        2/L3tint/{\faLightbulb}/L3/Counterfactual/{$P(E_{\scriptscriptstyle S{=}x} \mid S{=}s,\, E{=}e)$}/{``Had torque been raised at $t{=}20$\,ms, the arm would have stalled.''},%
        3/L4tint/{\faClipboardList}/L4/Decision Making/{$\pi^{*}\!=\!\arg\max_{\pi}\, U(\pi \mid E,S,C,K)$}/{``Given the observed torque spike, recommend a recovery procedure.''}%
    }{
        % Row background
        \fill[panelbg, rounded corners=3pt]
            (-0.15, \i*\rowsep - 0.02) rectangle (12.75, \i*\rowsep + 1.32);

        % Left accent bar
        \fill[\tint, rounded corners=2pt]
            (0, \i*\rowsep) rectangle (0.28, \i*\rowsep + 1.30);

        % Icon circle
        \node[
            circle, fill=\tint!15, draw=\tint, line width=0.7pt,
            minimum size=0.82cm, inner sep=0pt,
            font=\normalsize, text=\tint!80!darknavy
        ] (icon\i) at (0.92, \i*\rowsep + 0.65) {\icon};

        % Level label (bold) + subtitle
        \node[anchor=west, font=\sffamily\bfseries\small, text=\tint!80!darknavy]
            at (1.52, \i*\rowsep + 0.98) {\lbl\;\,$\cdot$\;\,\subtitle};

        % Formula
        \node[anchor=west, font=\sffamily\scriptsize, text=darknavy]
            at (1.52, \i*\rowsep + 0.62) {\formula};

        % Example question (italic)
        \node[anchor=west, font=\sffamily\tiny\itshape, text=figsteel,
              text width=10.5cm]
            at (1.52, \i*\rowsep + 0.28) {\example};
    }

    % Ascending arrow on the far right to indicate increasing reasoning complexity
    \draw[darknavy!50, line width=1.2pt, -{Stealth[length=4pt, width=3pt]}]
        (12.95, 0.15) -- (12.95, 3*\rowsep + 1.15);
    \node[anchor=south, font=\sffamily\tiny, text=darknavy!60, rotate=90]
        at (13.22, 2.0) {reasoning complexity};

    % Bracket on the left labelling L1--L3 as "Pearl's Ladder"
    \draw[decorate, decoration={brace, amplitude=4pt, mirror},
          darknavy!40, line width=0.6pt]
        (-0.40, -0.02) -- (-0.40, 2*\rowsep + 1.32);
    \node[anchor=east, font=\sffamily\tiny, text=darknavy!50, rotate=90]
        at (-0.72, 1.5*\rowsep + 0.30) {Pearl's Ladder};

    % Bracket for L4 labelled "FactoryBench"
    \draw[decorate, decoration={brace, amplitude=4pt, mirror},
          L4tint!70, line width=0.6pt]
        (-0.40, 3*\rowsep - 0.02) -- (-0.40, 3*\rowsep + 1.32);
    \node[anchor=east, font=\sffamily\tiny\bfseries, text=L4tint!80!darknavy, rotate=90]
        at (-0.72, 3*\rowsep + 0.65) {FactoryBench};

    \end{tikzpicture}
    \caption{Pearl's three rungs of causal reasoning (L1--L3) instantiated on robotic time-series data, plus the L4 decision-making layer introduced by FactoryBench. Each level demands progressively deeper reasoning: from passive state reading (\textbf{L1}), through predicting the effect of interventions (\textbf{L2}), to imagining alternative histories (\textbf{L3}), and finally synthesising goal-directed action sequences (\textbf{L4}). Color intensity encodes difficulty.}
    \label{fig:pearl_ladder}
\end{figure}
